\chapter{Conclusão}\label{cap:conclusao}

Este trabalho analisou e avaliou o \texttt{mcp-kicad-vscode}, o servidor
\emph{Model Context Protocol} de código aberto do projeto
\texttt{KiCAD-MCP-Server}, criado por \emph{mixelpixx}~\citep{kicadmcpserver},
voltado à automação do fluxo de projeto de placas de circuito impresso no
KiCad. O servidor expõe 233
ferramentas em 24 módulos funcionais, 173 delas descobríveis por busca em 16
categorias, além de 23 recursos dinâmicos de estado do projeto e quatro
\emph{prompts} de tarefas recorrentes. Integrações com JLCPCB, LCSC e DigiKey
e suporte ao autorouter Freerouting completam o ciclo de projeto até a
preparação de fabricação.

O objetivo geral foi atingido na medida em que o servidor foi analisado,
documentado e verificado de forma reprodutível, e o estudo de caso produziu
evidências de uso real. A avaliação combinou métricas
de implementação (68.819 linhas de código somando as duas camadas), execução
das suítes de teste (2.412 de 2.451 casos aprovados, com a suíte TypeScript
integralmente aprovada) e um estudo de caso, conduzido pelo autor, com uma
placa seguidora
de linhas de quatro camadas, com 97 \emph{footprints} e 83 redes, cuja
verificação de regras reportou 11 violações residuais e zero itens não
conectados.

Os objetivos específicos também foram cumpridos: a arquitetura e o protocolo
interno foram analisados e documentados (Capítulo~\ref{cap:arquitetura}); o
catálogo de ferramentas foi caracterizado e descrito
(Capítulo~\ref{cap:catalogo}); os mecanismos de descoberta e recursos foram
avaliados e medidos (Capítulo~\ref{cap:resultados}); as suítes de teste e a
integração contínua foram executadas; o estudo de caso, conduzido pelo autor,
produziu evidências concretas; e limitações e pendências foram registradas
(Capítulo~\ref{cap:discussao}).

As limitações do trabalho delimitam seu alcance com honestidade: a integração
IPC não foi exercitada no ambiente de validação, 60 ferramentas ainda não são
descobríveis, oito testes mostraram-se sensíveis ao ambiente, o projeto
estudado possui pendências de fabricação e não houve estudo com usuários. Cada
uma dessas limitações origina trabalho futuro diretamente acionável, o que
reforça o caráter incremental da contribuição.

Como consideração final, o trabalho demonstra que conectar assistentes de IA
a um fluxo de engenharia completo --- e não apenas a tarefas isoladas --- é
viável com um contrato padronizado, arquitetura em camadas e disciplina de
verificação. A adoção do MCP como espinha dorsal permite que o investimento
em integração seja compartilhado entre assistentes, e o modelo de catálogo
com descoberta oferece um caminho concreto para que agentes operem
ferramentas complexas de EDA com segurança e rastreabilidade.
